\chapter{Dataset and Preprocessing}
\label{ch:data}

\section{QASPER and the experimental population}
QASPER contains 5,049 information-seeking questions associated with 1,585
NLP research papers, with answers and supporting evidence annotations
\cite{dasigi2021qasper}. This structure makes it possible to evaluate
retrieval against evidence from a known source. It also introduces
practical complications: papers have different lengths, questions can
have multiple annotations, and some evidence refers to tables or figures
that do not have a directly recoverable text span.

The published dataset size is not the denominator used in every
repository experiment. The frozen supervised population contains 2,245
training questions from 845 papers and 924 validation questions from
277 papers. Those are the counts used for the main question-level
comparisons. They describe the preserved, eligible experimental records,
not a claim that every original question has usable text evidence for
every later method.

\begin{table}[tbp]
\centering
\caption{Populations used in the main experiments. The local training
rows are derived from the frozen question-level training population.}
\label{tab:populations}
\begin{tabularx}{\textwidth}{@{}Xrr@{}}
\toprule
Population & Questions & Papers\\
\midrule
Frozen question-level training & 2,245 & 845\\
Frozen validation retrieval & 924 & 277\\
Eligible local-router training & 2,101 & ---\\
Excluded from local training: no mappable text evidence & 144 & ---\\
\bottomrule
\end{tabularx}
\end{table}

The same paper must not appear on both sides of a grouped training
fold. The later five-fold procedures therefore group questions by
paper identifier. Validation remains outside those training folds.
However, it was used repeatedly across the development sequence,
including checkpoint selection in the Qwen experiments. This thesis
does not report a final test-set evaluation.

\section{Document construction and token boundaries}
\subsection{Flattening the source}
The ingestion code constructs a single paper text from the title,
abstract, section names, and section paragraphs. Nonempty components
are joined with double newlines. Retaining section names gives the
retriever access to signals such as an experimental or methodological
heading, while the single source string provides a consistent coordinate
system for chunk and evidence offsets.

This representation is a text view of a scientific paper, not a
reconstruction of its PDF layout. A table may be represented differently
from its appearance on the page, and figure-based evidence may not be
locatable as ordinary prose. The method should therefore be understood
as text evidence retrieval over the dataset's representation.

\subsection{Five aligned granularities}
The GPT-2 fast tokenizer is applied without special tokens and with
offset mappings. For a token sequence of length $L$, a size $g$ defines
successive ranges
\begin{equation}
[jg,\min((j+1)g,L)),\qquad
j=0,\ldots,\left\lceil L/g\right\rceil-1.
\label{eq:chunks}
\end{equation}
The final chunk may be shorter than $g$. The source text for a chunk is
sliced using the corresponding character offsets. Thus the stored
content can be traced to the flattened paper rather than reconstructed
from a potentially different token decoder.

Chunks do not overlap within a level. Because all sizes are multiples
of ten and double successively, boundaries align in token space across
levels. A complete 160-token root contains two 80-token children, four
40-token descendants, eight 20-token descendants, and sixteen 10-token
leaves. At document edges the tree can be partial. This alignment
supports the later hierarchy without implying that every boundary
coincides with a sentence or that every inter-token whitespace character
belongs to a stored span.

Token counts used for boundaries must also be distinguished from counts
after joining or normalising retrieved text. Joining separately sliced
chunks can change tokenisation at boundaries. For this reason, reported
retrieved token counts need not equal exactly five times the nominal
chunk size, even before considering short final chunks.

\section{Embedding and storage}
The project embeds questions, chunks, and evidence using
\path{text-embedding-3-small}, with 1,536 dimensions. Qdrant contains
three relevant collections: \texttt{PaperChunk},
\texttt{PaperQuestion}, and \texttt{PaperEvidence}. Payloads retain
identities, split information, chunk size, and available source offsets.
Granularity is a filter within the chunk collection; the five levels
should not be mistaken for five unrelated embedding models.

A read-only collection audit recorded 1,701,822 chunks, 4,526 questions,
and 9,522 evidence records. These are storage counts across the ingested
material. They are not the frozen training and validation sample sizes
in Table~\ref{tab:populations}. In particular, the presence of other
split records in storage does not imply that a test benchmark was
evaluated in the thesis.

The embedding pipeline batches requests, retries failures, and preserves
the association between input texts and returned vectors. Stable UUID5
identifiers are derived from document, granularity, and chunk index for
chunks; evidence identity incorporates its text hash. This allows
repeated ingestion to upsert the intended objects rather than silently
creating additional copies.

The Qdrant setup uses cosine distance, on-disk vector storage, HNSW,
and payload indexes for common filters. A historical audit recorded a
server/client version difference, with server 1.13.2 and client 1.16.1.
That observation belongs to the environment history; it is not evidence
that later frozen outputs are invalid. The repository also contains
legacy Weaviate-related code. The reported experiment pipeline uses
Qdrant, so those files are not an additional evaluated vector-store
baseline.

\section{Resumability and data integrity}
\subsection{From smoke tests to a frozen benchmark}
Early smoke tests were intended to validate execution, not estimate
retrieval performance. A February smoke artifact included an empty
output and a separate 15-record result corresponding to three questions
at five sizes. Neither is a full validation run. Treating these small
files as final results would confuse pipeline development with model
evaluation.

Ingestion was made resumable with seven tracked stages: one for each
chunk size, one for questions, and one for evidence. Checkpoint updates
use locking and atomic replacement. JSONL output writing also required
care: a bookkeeping problem in the locked writing/counting path was
addressed so that durable records and reported completion counts agreed.
These measures matter when embedding requests, retrieval, or training
are interrupted.

A June integrity audit covered 1,585 papers and identified four papers
with incomplete checkpoint state, six missing questions, sixteen missing
evidence records, and 1,345 evidence records without resolved offsets
out of 9,522. It also recorded 46 papers without highlighted evidence.
These are dated observations of the pipeline at that point. Subsequent
complete experiment artifacts, rather than a stale report of pending
work, determine the populations reported here.

The expected-regret phase provides a concrete example. Its utility
table initially covered 2,223 training questions. Recovering the missing
22 completed the frozen population of 2,245; it did not add a new
training dataset. The final comparison uses the completed target set.

\subsection{Text evidence and source offsets}
Two uses of evidence need different kinds of completeness. Joined
lexical F1 can compare the retrieved text with a supplied evidence string
even when that string has no exact source offset. Local gold-overlap
training, in contrast, needs to determine which source region intersects
the evidence. It cannot construct a meaningful span target from an
unlocatable annotation.

The local-router construction therefore separates complete, partial,
and unmappable evidence. Of the 2,245 training questions, 1,915 have
complete mapped evidence and 186 retain a usable partial mapping. The
remaining 144 are excluded from local training. Together the 2,101
eligible questions yield 4,081 gold-overlapping tree examples.
Annotations involving \texttt{FLOAT SELECTED} are an important source
of unmappable table or figure evidence.

Excluding unmappable questions is a defensible construction rule, but
it changes the training population. The local router may consequently
be less representative of questions whose support depends on visual
or tabular material. The thesis does not turn these exclusions into
zero-length labels or silently treat them as short-evidence cases.
All 924 frozen validation questions remain in the primary local
retrieval evaluation.

\section{Joined evidence-token evaluation}
\subsection{Normalisation and aggregation}
The versioned primary metric is \texttt{qasper-token-prf-v2}, with
lowercasing, ASCII punctuation removal, and whitespace collapsing. Reference
evidence strings are deduplicated under the metric's normalisation and
combined. Retrieved chunks are joined in their returned order. The
normalised strings are then tokenised using GPT-2. Each token is
decoded separately, surrounding whitespace is stripped, and empty
decoded tokens are discarded before counting. Thus the multiset
elements are these decoded token strings, not raw tokenizer IDs.

Let $a_t$ and $b_t$ be the multiplicities of token $t$ in the joined
retrieved text and reference evidence. The overlap is
\begin{equation}
O(q)=\sum_t\min(a_t,b_t).
\end{equation}
Precision, recall, and joined retrieval F1 are
\begin{equation}
P(q)=\frac{O(q)}{\sum_t a_t},\qquad
R(q)=\frac{O(q)}{\sum_t b_t},\qquad
\retrievalf(q)=\frac{2P(q)R(q)}{P(q)+R(q)}.
\label{eq:retrievalmetric}
\end{equation}
The implementation assigns zero when an empty denominator or zero
combined precision and recall prevents an ordinary calculation.
The reported mean is the arithmetic mean of per-question F1, not
the F1 computed from corpus-wide token totals and not the mean of
individual chunk F1 scores.

The multiset definition is important. Repeating a token in several
retrieved chunks increases the retrieved count, but its credit is
capped by the reference count. Redundant chunks can therefore lower
precision even when each looks relevant. Conversely, a matching token
can receive lexical credit without occurring at the same source
position as the annotated evidence. This is not a positional overlap
metric or a semantic entailment measure.

\subsection{Diagnostic similarity measures}
The baseline reports also store query-to-chunk and evidence-to-chunk
similarity summaries. Query similarity is available during retrieval.
Evidence similarity uses the reference evidence and is an offline
diagnostic. For a chunk, the latter can be summarised by its maximum
similarity to an evidence item and then aggregated across retrieved
chunks.

These measurements should not be substituted for joined F1. A passage
can have a high query similarity and omit much of the annotated evidence.
A long passage can resemble an evidence embedding while containing a
large amount of additional text. The experiments use these diagnostics
to inspect behaviour, while keeping the versioned lexical metric as
the primary retrieval outcome.

\section{Three targets called an oracle}
\label{sec:oracles}
\subsection{Retrieval-derived labels and the offline bound}
For each question, the fixed-size benchmark stores
\begin{equation}
U(q,g)=\retrievalf(R_g(q),E(q)),\qquad g\in\granularities.
\end{equation}
The best available utility is
\begin{equation}
U^*(q)=\max_{g\in\granularities}U(q,g).
\label{eq:oracle}
\end{equation}
Averaging $U^*(q)$ gives an offline bound over those five saved actions.
It is not an upper bound over every possible mixture of chunks or every
future retrieval architecture.

The early router uses a legacy label rule: maximise joined F1 within
$10^{-6}$, then use evidence similarity to break ties, and finally prefer
the smaller size. A later exact utility argmax with smaller-size
tie-breaking is close but not identical. On validation, their size
counts are respectively $(137,235,283,187,82)$ and
$(138,235,282,187,82)$. The thesis preserves that distinction rather
than silently reassigning the earlier model's targets.

\subsection{Whole-question evidence length}
The Qwen and global similarity-tree classification sequence uses a
different target. Evidence strings are stripped, exactly deduplicated,
sorted, joined with newlines, and tokenised. If their total length is
$\ell(q)$, the label is the nearest candidate size:
\begin{equation}
y_{\mathrm{len}}(q)=
\underset{g\in\granularities}{\arg\min}\ |\ell(q)-g|,
\label{eq:lengthlabel}
\end{equation}
with the smaller size chosen at a midpoint and extreme lengths clipped
to the available range. This target does not inspect which size
retrieves the best evidence.

\begin{table}[tbp]
\centering
\caption{Whole-question evidence-length labels. Counts sum to the
frozen training and validation populations.}
\label{tab:labels}
\begin{tabular}{@{}lrrrrr@{}}
\toprule
Split & 10 & 20 & 40 & 80 & 160\\
\midrule
Training & 55 & 267 & 586 & 687 & 650\\
Validation & 13 & 81 & 178 & 232 & 420\\
\bottomrule
\end{tabular}
\end{table}

The most common training label is 80, whereas the most common validation
label is 160. Predicting the training majority is a legitimate constant
baseline. Choosing the validation majority after inspecting validation
labels is only a descriptive reference, not a training-derived method.
This shift also helps explain why accuracy alone can conceal poor
coverage of smaller classes.

\subsection{Local evidence overlap}
The final tree-level target uses only the evidence intersecting a given
160-token root. Intersections are combined as a union so that overlapping
annotations do not count the same source region repeatedly. The
tokenised local evidence length is mapped to the nearest candidate
size with the same smaller-midpoint convention.

Its training unit is a question--tree pair, not a question. Its label
describes a local evidence span, not the total evidence in the paper.
Consequently, its class distribution and classification F1 cannot be
directly compared with the global evidence-length labels. The common
point of comparison is the final retrieval output on the same questions,
with the change in candidate-selection protocol stated explicitly.

\section{Reproducibility boundary}
The thesis draws numerical results from preserved reports, metric files,
and per-question outputs. It does not rerun embedding, GPU training,
or retrieval to create a more favourable result. Model revisions,
configuration fingerprints, source hashes, and pre-evaluation predictions
are retained where the corresponding phase recorded them. The generated
tables and figures are linked to those saved artifacts in
Appendix~\ref{app:provenance}.

This boundary keeps the writing task separate from a new experiment.
It also leaves a clear route for future work: a new run should produce
a new identifiable artifact, with its population, model state, and
metric version stated, rather than overwriting the result that supports
an earlier claim.
